\authorchapter{Ausgangslage und Zielsetzung}{Ereny Micail}
\label{ch:ausgangslage}

\section{Ausgangslage} 

In schulischen Teamprojekten läuft es meistens so: Aufgaben kommen in ein Tool,
Absprachen laufen über einen Messenger und Dateien liegen irgendwo anders.
Man wechselt ständig zwischen verschiedenen Apps hin und her, Infos gehen verloren
und keiner weiß mehr genau, wer gerade woran arbeitet.
Das nennt man \textit{Context Switching} -- und genau das war die Ausgangssituation
für \textit{EffSync}.

\bigskip

Wir wollten eine Lösung bauen, die Task-Management, Berechtigungen und
Echtzeit-Kommunikation in einem einzigen System zusammenbringt.
Der Nutzen: bessere Übersicht über den Arbeitsstand, einfachere Zusammenarbeit
zwischen Schülern und Lehrern und weniger organisatorischer Aufwand im Projektalltag.

\bigskip

Vor dem eigentlichen Projekt haben wir schon Grundlagen gelegt:
Architekturentscheidungen getroffen (Vue.js im Frontend, NestJS im Backend,
Prisma und PostgreSQL für die Datenbank), ein erstes Datenmodell erstellt und
erste Module prototypisch umgesetzt.
Dabei hat sich gezeigt, dass vor allem die konsistente Modellierung von
Task-Hierarchien, Berechtigungen und Echtzeit-Events entscheidend für die
Stabilität des Systems ist.


\authorsection{Auftrag}{ Ereny Micail}

Der Auftrag war klar: eine webbasierte Plattform bauen, auf der mehrere Nutzer
gemeinsam Aufgaben verwalten und direkt miteinander kommunizieren können.
Aufgaben sollen strukturiert erfasst, zugewiesen und verfolgt werden können --
inklusive Unteraufgaben und klaren Verantwortlichkeiten.
Absprachen sollen direkt im System stattfinden, damit der fachliche Kontext
immer mit den zugehörigen Aufgaben verknüpft bleibt.

\bigskip

Dazu kommt ein Berechtigungskonzept: Nur wer dazu berechtigt ist, darf
sicherheitsrelevante Aktionen ausführen.
Die technische Umsetzung soll nachvollziehbar dokumentiert, durch Tests
abgesichert und wissenschaftlich begründet sein.


\authorsection{Projektziele}{ Ereny Micail}

Für das aktuelle Semester haben wir folgende SMART-Ziele definiert:

\begin{enumerate}
    \item \textbf{Task-Hierarchie produktiv nutzbar umsetzen:}
          Bis zum Semesterende können Benutzer Tasks inklusive Subtasks in
          mindestens drei Hierarchieebenen erstellen, bearbeiten und abrufen.
          Zirkelbezüge werden serverseitig erkannt und mit HTTP~400 abgewiesen.

    \item \textbf{Berechtigungslogik für kritische Aktionen durchsetzen:}
          Bis zum Semesterende sind für Task-Mitgliederverwaltung und
          Chatraum-Verwaltung rollen- bzw.\ kontextbasierte Prüfungen
          implementiert, sodass unautorisierte Änderungsversuche reproduzierbar
          mit HTTP~403 blockiert werden.

    \item \textbf{Echtzeitkommunikation integrieren:}
          Bis zum Semesterende ist ein funktionsfähiger Live-Chat mit
          WebSocket-basierter Zustellung umgesetzt, bei dem Nachrichten in
          privaten und gruppenbezogenen Räumen ohne manuelles Neuladen auf
          verbundenen Clients erscheinen.

    \item \textbf{Qualitätssicherung über automatisierte Tests etablieren:}
          Bis zum Semesterende existiert eine lauffähige Testbasis aus Unit- und
          E2E-Tests (\texttt{npm run test}, \texttt{npm run test:e2e}), die
          zentrale Kernabläufe (Auth, Throttling, Mitgliederverwaltung,
          Berechtigungen) reproduzierbar validiert.
\end{enumerate}

\subsection*{Abgrenzung (Nichtziele)}

Nicht Teil dieser Etappe sind native Mobile-Apps, Mandantenfähigkeit für mehrere
Schulen oder Organisationen, horizontale Hochskalierung für sehr große Nutzerzahlen
sowie vollständige Offline-Funktionalität.


\authorsection{Systemgrenze}{ Ereny Micail}

\textit{EffSync} besteht aus dem Vue.js-Frontend, dem NestJS-Backend, den REST-
und WebSocket-Schnittstellen sowie der Datenbank über Prisma und PostgreSQL.
Externe Systeme wie E-Mail-Dienste oder Drittplattformen werden nur über klar
definierte Schnittstellen angebunden und sind nicht Teil der Kernlogik.
Nutzer arbeiten ausschließlich über das Frontend -- Persistenz, Validierung und
Autorisierung laufen im Backend.


\authorsection{Vorgehen}{ Ereny Micail}

Wir setzen das Projekt iterativ um.
Zuerst haben wir Datenmodell und Kernlogik (Tasks, Auth, Berechtigungen)
stabilisiert, danach kamen Echtzeit-Funktionen und Usability-Verbesserungen dazu.
Die Abhängigkeiten zwischen den Zielen haben wir dabei bewusst berücksichtigt:
Eine sichere Mitgliederverwaltung braucht ein konsistentes Rollenmodell,
und Echtzeit-Updates brauchen valide Zustandsänderungen im Backend.
Risiken haben wir durch frühe Prototypen, laufende Reviews und automatisierte
Tests reduziert.
